% math-article-template: 1
% Replace every instructional sentence before delivery.
% The same source is used for review and final publication.
% Set \draftpaperfalse after review to hide the draft date/version and line numbers.
\newif\ifdraftpaper
\draftpapertrue

\documentclass[11pt]{article}

\usepackage[utf8]{inputenc}
\usepackage[a4paper,margin=30mm]{geometry}
\usepackage{amsmath,amssymb,amsthm,mathtools}
\usepackage{microtype}
\usepackage{xcolor}
\ifdraftpaper
  \usepackage[mathlines]{lineno}
\fi
\usepackage[
  colorlinks=true,
  linkcolor=blue!55!black,
  citecolor=green!35!black,
  urlcolor=blue!65!black
]{hyperref}

\newtheorem{theorem}{Theorem}[section]
\newtheorem{proposition}[theorem]{Proposition}
\newtheorem{lemma}[theorem]{Lemma}
\newtheorem{corollary}[theorem]{Corollary}
\theoremstyle{definition}
\newtheorem{definition}[theorem]{Definition}
\newtheorem{remark}[theorem]{Remark}

\newcommand{\paperversion}{0.1}
\title{Replace with a mathematical title}
\author{Replace with the author name}
\ifdraftpaper
  \date{Draft \paperversion\ --- \today}
\else
  \date{}
\fi

\begin{document}
\maketitle
\ifdraftpaper
  \linenumbers
\fi

\begin{abstract}
State the main result, principal hypotheses, method, and significance in a
self-contained paragraph. Use only terminology defined in the paper or
standard for the intended audience.
\end{abstract}

\section{Introduction}

Explain the mathematical problem and its context without referring to project
history or drafting chronology. State what is new, why the result matters, and
the precise scope of the paper.

\begin{theorem}[Main theorem]\label{thm:main}
State the exact main theorem with all hypotheses and quantifiers.
\end{theorem}

Give a compact proof strategy that names the actual mathematical mechanisms.
End the introduction with a short description of the paper's logical
organization.

\section{Conventions and notation}

Fix the ambient categories, base objects, variance, finiteness and completeness
hypotheses, and nonstandard notation. Define every abbreviation before use and
use each symbol consistently.

\section{Preliminaries and external inputs}

State only the definitions and results genuinely needed later. For each
external result, give its exact usable statement, a precise citation, and the
bridge from its hypotheses to the objects of this paper.

\begin{proposition}[Descriptive name of an external input]\label{prop:external}
State the complete consequence used in the main argument.
\end{proposition}

\section{Main argument}

Introduce each lemma by explaining its purpose and connection to the main
theorem. Avoid new terminology for a construction used only once.

\begin{lemma}[Descriptive structural lemma]\label{lem:structural}
State the complete structural lemma.
\end{lemma}

\begin{proof}
Give the proof, with every map, comparison, and change of ambient category made
explicit when it carries mathematical content.
\end{proof}

\section{Proof of the main theorem}

\begin{proof}[Proof of Theorem~\ref{thm:main}]
State the consequences of Proposition~\ref{prop:external} and
Lemma~\ref{lem:structural} being applied, verify their hypotheses in the
present setting, complete the argument, and identify the conclusion of
Theorem~\ref{thm:main} explicitly.
\end{proof}

% Include only references actually used. Give exact theorem or definition
% locators in the prose where each source enters.
% \begin{thebibliography}{9}
% \bibitem{source-key} Author, \emph{Title}, version, publication data.
% \end{thebibliography}

\end{document}
